% No modificar estas líneas de código, por favor dirigirse a RESUMEN y ABSTRACT
\newpage

%------------------------
%
%       RESUMEN 
%
%------------------------

\section*{RESUMEN}

El presente proyecto de grado describe el diseño, implementación y validación experimental de un sistema embebido periférico y no invasivo de monitoreo de condición de herramientas (\textit{Tool Condition Monitoring}, TCM) para tornos de Control Numérico Computarizado (CNC) en la empresa metalmecánica Maestranza San Carlos (MAINSA). La investigación responde a la necesidad de superar la supervisión empírica visual tradicional en taller, la cual genera roturas intempestivas de insertos, pérdidas de precisión dimensional, acabados superficiales deficientes y paradas imprevistas de producción. La arquitectura mecatrónica propuesta integra la adquisición no invasiva de señales de corriente del motor del husillo mediante un transformador de corriente SCT013 y vibraciones mecánicas a través de un acelerómetro, acoplados a una etapa de acondicionamiento analógico para conversión a niveles proporcionales al valor eficaz (RMS). El procesamiento y la extracción de características se ejecutan en tiempo real sobre un microcontrolador STM32, el cual aloja un modelo ligero de Inteligencia Artificial en el borde (\textit{Edge AI}) optimizado para clasificar los niveles de desgaste y predecir eventos de falla con una latencia inferior a 20 ms. Paralelamente, el sistema cuenta con conectividad IoT para la transmisión inalámbrica de métricas operativas y alertas tempranas a un panel de supervisión remoto. La validación en planta demostró una reducción sustancial en piezas defectuosas y tiempos muertos por retrabajo, consolidando una solución de mantenimiento predictivo costo-efectiva, modular y de bajo impacto para maquinaria industrial heredada.

\textbf{Palabras clave:} Monitoreo de condición de herramientas, torno CNC, sensor SCT013, acelerometría, Edge AI, STM32, Internet de las Cosas (IoT), mantenimiento predictivo, MAINSA.

\newpage

%------------------------
%
%       ABSTRACT
%
%------------------------

\section*{ABSTRACT}

This graduation project describes the design, implementation, and experimental validation of a peripheral, non-invasive embedded system for Tool Condition Monitoring (TCM) applied to Computer Numerical Control (CNC) lathes at the Maestranza San Carlos (MAINSA) metalworking facility. This research addresses the operational limitations of traditional empirical visual inspection in workshops, which frequently results in sudden cutting insert breakages, dimensional inaccuracies, substandard surface finishes, and unscheduled production downtime. The proposed mechatronic architecture incorporates non-invasive acquisition of spindle motor current signals via an SCT013 split-core current transformer and mechanical vibrations via an accelerometer, interfaced with an analog conditioning stage for root mean square (RMS) conversion. Signal processing and feature extraction are executed in real time on an STM32 microcontroller hosting an optimized Edge Artificial Intelligence (Edge AI) lightweight model to classify tool wear severity and anticipate critical tool failures with an inference latency under 20 ms. Concurrently, the system features IoT connectivity for wireless telemetry and early warning alerts through a remote supervisory dashboard. Experimental validation in real workshop operations demonstrated a substantial reduction in discarded workpieces and rework downtime, delivering a cost-effective, modular, and low-impact predictive maintenance solution for legacy industrial machine tools.

\textbf{Keywords:} Tool condition monitoring, CNC lathe, SCT013 sensor, accelerometry, Edge AI, STM32, Internet of Things (IoT), predictive maintenance, MAINSA.